\section{引言}

大语言模型（LLM）的出现极大提升了机器在自然语言理解与生成方面的能力，而近来的多模态大语言模型进一步把这种能力扩展到了视觉感知与动作规划~\citep{Bai2025Qwen3VL, comanici2025gemini_gemini25, gpt5, claudeopus45, guo2025seed1_seed15}。这为 GUI 自动化奠定了基础：模型不仅能够理解复杂界面，还能够推理出合适的交互步骤，从而成为真正可用的图形界面智能体。

围绕 GUI 智能体，研究社区已经在数据采集、模型训练和任务评测上取得显著进展。已有工作探索了 GUI 交互数据的获取方法~\citep{gao2024mobileviews, rawles2023androidinthewild, liu2025scalecua, wang2025opencua}，并提出了多种具备实际能力的 GUI 模型~\citep{qin2025ui_uitars1, wang2025ui_uitars2, ye2025mobile_mobileagentv3, zeng2025uitron}。然而，一个关键瓶颈始终存在：\emph{如何高效获得高质量的轨迹和知识数据，以持续提升 GUI 智能体在目标领域中的能力？} 传统人工标注不仅昂贵，而且容易受主观因素影响，难以支撑大规模迭代。

为解决这一问题，本文提出以 \textbf{Calibrated Step Reward System}（CSRS）为核心的自进化训练流水线。与依赖模型自反思、后验过滤或单轮可验证任务的数据生成方法不同~\citep{zelikman2022star, madaan2023self_selfrefine, jung2023impossibledistill, gulcehre2023reinforced_rest}，CSRS 面向多步 GUI 执行场景，把模型生成的步骤级推理锚定到轨迹级结果信号上，这些结果信号可以来自自动验证脚本，也可以来自人工标注。通过这种方式，CSRS 能在可扩展的前提下显著提高轨迹标注质量，并将生成轨迹进一步提炼为训练所需的知识数据。

在这套流水线上，我们训练了 \textbf{Step-GUI} 系列模型（4B 和 8B），其骨干来自 Qwen3-VL。Step-GUI-8B 在多个 GUI 基准上达到领先水平：AndroidWorld 80.2\%，OSWorld 48.5\%，ScreenShot-Pro 62.6\%，并在本文提出的 AndroidDaily 上获得静态任务 89.91\%、端到端任务 52.50\% 的结果。与此同时，4B 版本已经具备较强的本地部署能力，可运行在消费级硬件上，适合隐私敏感场景。

当 GUI 智能体逐渐具备更强的理解与执行能力后，又会出现两个更现实的部署问题：一是不同设备和平台之间缺乏统一的通信接口，二是屏幕截图和设备状态往往包含大量敏感信息。为此，我们提出 \textbf{GUI-MCP}，即面向 GUI 自动化的图形界面模型上下文协议。它通过分层双层设计统一设备控制接口，在低层提供原子级操作，在高层支持把任务直接委派给本地 GUI 专家模型，从而让主模型聚焦高层规划，同时把实际界面执行尽可能留在本地设备上，以实现更好的隐私保护。

除了训练和部署，真实可用的 GUI 智能体还需要更贴近现实的评测。现有工作要么偏重静态界面 grounding~\citep{cheng2024seeclick, wu2024atlas, li2025screenspotpro, li2025autogui, burns2022dataset_motif, rastogi2021uibert_refexp, zhang2024llamatouch, liu2024visualwebbench_vwb}，要么虽然支持交互式执行~\citep{rawles2024androidworld, xie2024osworld}，但对高频中文移动应用场景覆盖不足。针对这一空白，我们提出 \textbf{AndroidDaily}。该基准不追求最大化应用数量，而是聚焦出行、购物、社交、娱乐和本地服务等高频真实场景，并采用静态动作预测与端到端执行相结合的双层评测方式，更能反映 GUI 智能体在日常数字生活中的实际可用性。

总体来看，本文的贡献体现在四个方面：
\begin{itemize}
\item 提出自进化训练流水线与 CSRS，以大幅降低标注成本并持续提升模型能力。
\item 训练出 Step-GUI 系列模型，在多个 GUI 基准上达到领先或有竞争力的结果，同时保留较强通用能力。
\item 提出 GUI-MCP，为 GUI 自动化提供统一、可扩展且隐私友好的协议接口。
\item 提出 AndroidDaily，建立更贴近日常移动使用场景的 GUI 评测基准。
\end{itemize}

这些工作共同覆盖了 GUI 智能体从训练数据、模型构建、部署接口到真实评测的完整链条，为实用型 GUI 智能体奠定了基础。
